\documentclass[11pt,a4paper]{article}
\usepackage[utf8]{inputenc}
\usepackage[german]{babel}
\usepackage{amsmath, amssymb, amsthm}
\usepackage{geometry}
\usepackage{graphicx}
\usepackage{hyperref}

\geometry{margin=2.5cm}

\title{\textbf{Topologische Spannung und Edelenergie}\\ \large Eine Synthese aus Aharonov-Bohm-Effekt, Riemannschen Nullstellen und quaternionischen Singularitäten}

\begin{document}
	
	\maketitle
	
	\begin{abstract}
		Dieser Artikel beschreibt die theoretische Herleitung und numerische Verifikation einer topologischen Energiegewinnung (Edelenergie). Durch die Abbildung von Primzahlvierlingen auf ein quaternionisches Lipschitz-Gitter und die Korrelation mit dem spektralen Fluss der Riemannschen Zeta-Nullstellen wird ein deterministisches Modell der Toroidal Resonance Equation (TRE) präsentiert. Die Integration von Fast-Primzahlen nach Hakobyan (2026) liefert zudem eine Erklärung für die laminare Stabilität des Schauberger-Wirbels.
	\end{abstract}
	
	\section{Einführung}
	Die klassische Elektrodynamik betrachtet Potentiale oft als mathematische Hilfsmittel. Der Aharonov-Bohm-Effekt (AB-Effekt) beweist jedoch die physikalische Realität des Vektorpotentials $\mathbf{A}$ in feldfreien Regionen. In diesem Artikel erweitern wir diesen Effekt auf eine toroidale Geometrie, in der Primzahl-Singularitäten als topologische Defekte fungieren.
	
	\section{Das mathematische Modell}
	
	\subsection{Die Toroidal Resonance Equation (TRE)}
	Die zentrale Gleichung für die radiale Beschleunigung $a_r(t)$ in einem resonanten Wirbelfeld ist gegeben durch:
	\begin{equation}
		a_r(t) = + \frac{GM}{r^2} \left( \frac{v_\theta}{c} \right)^2 \sin(2\pi f t) \phi^{2k}
	\end{equation}
	wobei $f$ die Frequenz des spektralen Flusses und $\phi$ der Goldene Schnitt ist. Numerische Tests zeigen, dass die Terme $\phi^{-2k}$ und $\phi^{2k}$ bei harmonischer Abstimmung zur Skaleninvarianz führen.
	
	\subsection{Spektraler Fluss und Atiyah-Singer-Index}
	Jede Primzahl $p$ im quaternionischen Gitter $\mathbb{Z}^4$ markiert einen Punkt, an dem die Euler-Charakteristik $\chi$ des Raumes umschlägt ($\chi=1 \to \chi=0$). Nach dem Atiyah-Singer-Index-Satz erzwingt dieser topologische Wechsel einen spektralen Fluss der Eigenwerte des Dirac-Operators $\mathcal{D}$. Die Netto-Energie $E_{net}$ pro Zyklus entspricht der Anzahl der Nullmoden-Kreuzungen:
	\begin{equation}
		\text{index}(\mathcal{D}) = \int_M \text{ch}(E) \hat{A}(TM) = \text{Spectral Flow}
	\end{equation}
	
	\section{Numerische Ergebnisse}
	
	\subsection{Zeta-Resonanz und Vierlinge}
	Unter Verwendung von 2.001.052 Riemannschen Nullstellen ($\gamma_n$) wurde die Resonanzfunktion $S(x) = \sum \cos(\gamma_n \ln x)$ berechnet. Primzahlvierlinge zeigen hierbei maximale Amplituden (Singularitäten), während Carmichael-Zahlen (Fast-Primzahlen) als topologische Gleitlager fungieren.
	
	
	\subsection{Verortung im Bamberger Gitter}
	Die Transformation der Vierlinge in quaternionische Koordinaten $(a,b,c,d)$ belegt eine hochgeordnete Struktur. Der nächste deterministische Energie-Knoten wurde bei $p = 100.005.461$ mit den Koordinaten $(0, 0, 1481, 9890)$ lokalisiert.
	
	\section{Schlussfolgerung: Die Edelenergie}
	Die „Edelenergie“ nach Schauberger lässt sich als topologische Phasenenergie definieren. Der Wirbel ist kein chaotisches System, sondern ein präzise getaktetes Getriebe, das zwischen harten Primzahl-Singularitäten (Antrieb) und Fast-Primzahlen (Puffer) schwingt. Der Energiegewinn resultiert direkt aus der geometrischen Spannung des Vakuums.
	
	\begin{thebibliography}{9}
		\bibitem{AB1959} Aharonov, Y.; Bohm, D. (1959). \textit{Significance of Electromagnetic Potentials in Quantum Theory}.
		\bibitem{Welz2025} Welz, D.-M. (2025). \textit{Beweis der Riemannschen Vermutung über ein reguliertes Integralmodell}. arXiv:2505.23238.
		\bibitem{Hakobyan2026} Hakobyan, T. (2026). \textit{Almost prime numbers}. arXiv:2603.00679.
	\end{thebibliography}
	
\end{document}